% Exact endpoint-image reduction, owned by cycles. No preamble.
\paragraph{C.1n (Proved: the endpoint image criterion is exactly the
universal standard conjecture $D$).}
We first work over $\mathbb C$ with rational coefficients. Fix integers
$n\ge0$ and $0\le r\le n$. Let $\mathcal I_{n,r}$ mean the
following assertion: for every smooth projective connected
$n$-fold $X$, every smooth projective connected $d$-fold $Y$,
every $0\le j\le d$, and every correspondence
\[
 \Gamma\in\mathrm{CH}^{d+r-j}(Y\times X)_{\mathbb Q},
\]
an element $z\in A^r_{\mathrm{alg}}(X)$ belongs to the image
of $\Gamma_*:A^j_{\mathrm{alg}}(Y)\to A^r_{\mathrm{alg}}(X)$
if and only if its cycle class belongs to the image of
\[
 \Gamma_*:H^{2j}(Y,\mathbb Q(j))
                      \longrightarrow H^{2r}(X,\mathbb Q(r)).
\]
These are the source and degree-shift quantifiers of the image
criterion on manuscript page $14{:}46$: its source degree is
arbitrary, and the correspondence need not have degree zero.
Write $D_{\mathbb C}$ for equality of numerical and Betti
homological equivalence on rational cycles of every smooth
projective complex variety. Then
\begin{equation}\label{eq:endpoint-image-D}
 \mathcal I_{n,r}\quad\Longleftrightarrow\quad D_{\mathbb C}
       \qquad\text{if }r\in\{0,1,n\}.
\end{equation}
For every $2\le r\le n-1$, the assertion $\mathcal I_{n,r}$
is false by C.1k. If we restrict the target to
$X=\mathbb P^n$, its image criterion is equivalent to
$D_{\mathbb C}$ for \emph{every} $0\le r\le n$.
More precisely, for a \emph{fixed} $X$ and $r$, let
$\mathcal I(X,r)$ require the displayed criterion for every
source $Y$, source degree $j$, and correspondence $\Gamma$.
Then
\begin{equation}\label{eq:fixed-target-image-D}
 \mathcal I(X,r)\quad\Longleftrightarrow\quad
             \bigl(D_{\mathbb C}\ \text{and}\ G^r(X)=0\bigr).
\end{equation}

\emph{Necessity, including the restricted projective-space target.}
Suppose that $D_{\mathbb C}$ fails. Then some smooth projective
$d$-fold $Y$ has a numerically trivial rational cycle
$a\in\mathrm{CH}^s(Y)_{\mathbb Q}$ with
$\mathrm{cl}(a)\ne0$. Put $j=d-s$, let $h$ be the hyperplane
class on $\mathbb P^n$, and set
\[
 \Gamma=a\times h^r\in
       \mathrm{CH}^{s+r}(Y\times\mathbb P^n)_{\mathbb Q}.
\]
The codimension is $s+r=d+r-j$, as required. The action on
cycles is
\[
 \Gamma_*b=\deg(a\cdot b)h^r=0
       \qquad\bigl(b\in A^j_{\mathrm{alg}}(Y)\bigr),
\]
by numerical triviality. In contrast, Poincar\'e duality supplies
$v\in H^{2d-2s}(Y,\mathbb Q(d-s))$ with
$\int_Y\mathrm{cl}(a)v=1$. Hence
$\Gamma_*v=\mathrm{cl}(h^r)$, although the nonzero cycle class
$h^r\in A^r_{\mathrm{alg}}(\mathbb P^n)$ is not in the cycle
image. This contradicts the proposed criterion. Notice that
we neither need nor assert that this cohomological preimage
$v$ is algebraic: the criterion itself would force an
algebraic preimage, which numerical triviality excludes.
The same construction works for any fixed target $X$: take
$h$ ample on $X$; its power $h^r$ has nonzero cohomology
class, since $\int_Xh^n>0$. Thus $\mathcal I(X,r)$ also
implies $D_{\mathbb C}$. Taking $\Gamma=0$ shows separately
that $\mathcal I(X,r)$ implies $G^r(X)=0$.

\emph{Sufficiency.}
Assume $D_{\mathbb C}$. By \citet[Theorem~1, pp.~448--449]
{Jannsen1992}, the category of rational numerical motives
is semisimple abelian. Betti realization on this category is
well defined and faithful: $D_{\mathbb C}$ applies on every
product, and K\"unneth together with Poincar\'e duality
identifies a correspondence's total cohomology class with
its action on all cohomological degrees. Here K\"unneth is
only the cohomological decomposition; we do not assume
algebraic K\"unneth projectors.

We use the convention
\[
 \begin{split}
 \operatorname{Hom}(\mathbf1,h(X)(r))
       &=A^r_{\mathrm{num}}(X),\\
 \operatorname{Hom}(h(Y)(j),h(X)(r))
       &=A^{d+r-j}_{\mathrm{num}}(Y\times X).
 \end{split}
\]
Thus $\Gamma$ induces a morphism
$f:M=h(Y)(j)\to N=h(X)(r)$, and the numerical class of
$z$ is a morphism $\beta:\mathbf1\to N$. Semisimplicity
gives $g:N\to M$ for which $p=fg$ is a projector onto
$\operatorname{im}(f)$ and $pf=f$. If $\mathrm{cl}(z)$ is
in the stated cohomological image, then
$H((1-p)\beta)=0$. Indeed realization in degree zero of
$h(Y)(j)$ and $h(X)(r)$ is precisely the pair of groups in
the displayed image criterion. Faithfulness gives
$(1-p)\beta=0$, so
\[
                     \beta=f(g\beta).
\]
Choose a rational cycle $b\in\mathrm{CH}^j(Y)_{\mathbb Q}$
representing $g\beta$. We have proved that $\Gamma_*b$ and
$z$ are numerically equivalent. In target codimensions
$r=0,n$, numerical and algebraic equivalence with rational
coefficients are both determined by degree. In codimension
$r=1$, the same equality follows from the N\'eron--Severi
theorem and the divisor intersection pairing, as recalled
in C.1k. Therefore $\Gamma_*b=z$ in
$A^r_{\mathrm{alg}}(X)$ in every asserted endpoint case.
The converse implication is always given by functoriality
of the cycle class. For $X=\mathbb P^n$, the map
$A^r_{\mathrm{alg}}(X)\to A^r_{\mathrm{num}}(X)$ is an
isomorphism in every codimension, so the same proof gives
the final restricted-target assertion.
For the fixed-target assertion, $D_{\mathbb C}$ identifies
numerical and homological equivalence, and $G^r(X)=0$
identifies homological and algebraic equivalence in the
specified target group. The same numerical preimage therefore
proves \eqref{eq:fixed-target-image-D}.

These two equivalences are not restricted to complex geometry.
Let $k$ be any algebraically closed field and let $H$ be a
fixed Weil cohomology with characteristic-zero coefficient
field $F$. Keep rational cycle coefficients, use $F$ in the
cohomological image, and define $D_{k,H}$ by equality of
numerical and $H$-homological equivalence for rational cycles
on every smooth projective $k$-variety. The identical argument
proves
\[
 \mathcal I_{n,r}^{k,H}\ \Longleftrightarrow\ D_{k,H}
         \qquad(r\in\{0,1,n\}),
\]
\[
 \mathcal I^{k,H}(X,r)\ \Longleftrightarrow\
       \bigl(D_{k,H}\ \text{and}\ G_H^r(X)=0\bigr),
\]
where $G_H^r(X)$ is the kernel of the rational-cycle class map
into $H^{2r}(X)(r)$. The source functional is evaluated using
Poincar\'e duality over $F$; the numerical-motive category
and its image splitting still have rational coefficients.
Numerical and algebraic equivalence agree rationally for
divisors over any algebraically closed field, and degree
handles codimensions zero and top. The intermediate-codimension
universal counterexamples cited above remain assertions over
$\mathbb C$; this extension does not transfer their
nonvanishing to positive characteristic.

Equations~\eqref{eq:endpoint-image-D}--\eqref{eq:fixed-target-image-D}
isolate an exact remaining
standard-conjecture gate. It does not prove $D_{\mathbb C}$
and does not replace homological equivalence by algebraic
equivalence in intermediate target codimension. It also
explains why the absence of a target Griffiths group by
itself was insufficient for the image argument in C.1k:
the unrestricted source can still detect a failure of $D$.
